\section{Анализ предметной области}

\subsection{Продукционные системы Маркова: история и значение}

Продукционные системы Маркова (или нормальные алгоритмические системы) были предложены Андреем Андреевичем Марковым в 1950-х годах и стали одной из первых формальных моделей вычислений, эквивалентных по вычислительной мощности машине Тьюринга. Их появление стало важным этапом в развитии теории алгоритмов и формальных языков, заложив основы для изучения автоматов, грамматик и систем подстановок.

Модель Маркова используется для формализации процессов преобразования строк, доказательства вычислимости функций и анализа алгоритмических процедур. Благодаря простоте формулировки и универсальности принципа подстановок, продукционные системы находят применение в учебных целях, при разработке интерпретаторов формальных языков, систем логического вывода и генераторов текстов.

\subsection{Формальное определение и структура системы Маркова}

Нормальная алгоритмическая система Маркова (НСМ) определяется как упорядоченная совокупность:
\[
M = (A, P, w_0),
\]
где:
\begin{itemize}
	\item \(A\) — конечный алфавит символов;
	\item \(P\) — упорядоченный набор правил подстановок (продукций) вида \( \alpha_i \rightarrow \beta_i \) или \( \alpha_i \rightarrow \beta_i. \), где точка указывает на завершающее правило;
	\item \(w_0\) — начальное слово (входная строка).
\end{itemize}

Каждое правило читается следующим образом: если в текущем слове встречается подстрока \( \alpha_i \), то она заменяется на \( \beta_i \). Применяется первое по порядку правило, в котором найдено вхождение левой части. Если правило завершающее (с точкой), выполнение алгоритма останавливается. В противном случае процесс продолжается до тех пор, пока возможно применение хотя бы одного правила.

\subsection{Механизм функционирования и пример работы}

Алгоритм работы системы Маркова включает следующие этапы:
\begin{enumerate}
	\item Задание исходного слова \(w_0\);
	\item Последовательный просмотр списка правил \(P\);
	\item Поиск первой подстроки, совпадающей с левой частью правила;
	\item Замена найденной подстроки на правую часть;
	\item Проверка наличия точки в правиле: если она есть — завершение работы;
	\item Возврат к шагу 2 при возможности дальнейших замен.
\end{enumerate}

\textbf{Пример.} Пусть задана система:
\[
\begin{aligned}
	a &\rightarrow ab,\\
	b &\rightarrow c.,\\
	a &\rightarrow d.
\end{aligned}
\]
и начальное слово \(a\). Тогда последовательность преобразований:
\[
a \Rightarrow ab \Rightarrow ac,
\]
и алгоритм завершает работу, так как правило \(b \rightarrow c.\) является завершающим.

\subsection{Цели и задачи моделирования системы Маркова}

Разработка программной реализации системы Маркова направлена на моделирование процесса пошагового применения правил и визуализацию изменений строки во времени.  
Такой симулятор должен предоставлять следующие возможности:
\begin{itemize}
	\item задание входной строки и набора правил подстановок;
	\item пошаговое выполнение алгоритма с отображением каждого преобразования;
	\item автоматический режим выполнения с регулировкой скорости;
	\item сохранение журнала (логов) выполнения и статистики применения правил;
	\item визуальное отображение текущего состояния строки и номера применяемого правила;
	\item загрузку и сохранение конфигураций из текстовых файлов.
\end{itemize}

\subsection{Существующие решения и их ограничения}

Существуют различные онлайн-симуляторы и учебные программы, моделирующие систему Маркова, однако большинство из них не предоставляет возможности сохранять результаты работы, собирать статистику или пошагово демонстрировать процесс с визуальной анимацией. Кроме того, нередко такие решения не поддерживают локальное хранение правил и не позволяют интегрировать пользовательские конфигурации.

Использование языка \texttt{Python} в сочетании с библиотеками \texttt{tkinter} и \texttt{customtkinter} позволяет реализовать современный интерфейс, поддерживающий как ручной ввод данных, так и загрузку правил из файлов. Встроенная система логирования облегчает анализ хода выполнения, а визуальная анимация помогает наглядно изучить логику работы алгоритма.

\subsection{Выводы по предметной области}

Продукционные системы Маркова представляют собой простую, но мощную формальную модель вычислений. Их моделирование требует реализации чёткой последовательности применения правил, механизма остановки и средств визуализации преобразований. Разработка симулятора на Python с графическим интерфейсом обеспечивает наглядность, интерактивность и возможность анализа статистики выполнения, что делает его полезным инструментом для обучения и исследования формальных алгоритмов.
